\chapter{Civilproces, straffeproces og bevis}

\chapterlead{Procesretten er retten til at få en konflikt behandlet på en bestemt måde. Den afgør, hvem der kan sagsøge, hvad retten kan tage stilling til, hvordan bevis føres, og hvordan en afgørelse kan prøves og fuldbyrdes.}

\section{Materiel ret og procesret}

Materiel ret siger, hvad der gælder mellem parter: om en aftale er bindende, om en person er erstatningsansvarlig, eller om en handling er strafbar. Procesret siger, hvordan spørgsmålet behandles af en domstol eller et andet organ.

En person kan have et materielt krav uden at kunne få medhold i en sag, hvis kravet er forældet, dokumentationen er utilstrækkelig, den forkerte part er sagsøgt, eller sagen ikke er rejst korrekt. Proces er derfor ikke et teknisk tillæg; den former den praktiske adgang til ret.

Retsplejeloven indeholder de centrale regler om retternes organisation, civile sager, straffesager, bevis, appel og fuldbyrdelse \citep{retsplejeloven}.

\section{Civil sag}

En civil sag begynder typisk med, at sagsøgeren fremsætter et krav mod sagsøgte. Sagen skal identificere parterne, påstanden, de faktiske omstændigheder og den juridiske begrundelse. Retten tager som udgangspunkt stilling til det, parterne har bragt ind i sagen.

Forløbet kan indeholde:
\begin{enumerate}
  \item stævning og svarskrift,
  \item forberedelse, hvor krav og beviser afgrænses,
  \item eventuel retsmægling eller forlig,
  \item hovedforhandling med forklaringer og dokumenter,
  \item dom, og
  \item eventuel appel og fuldbyrdelse.
\end{enumerate}

\section{Straffesag}

I en straffesag er det staten ved anklagemyndigheden, der rejser sagen mod den tiltalte. Politiets efterforskning og anklagemyndighedens beslutning om tiltale er adskilt fra domstolens afgørelse. Den tiltalte har ret til forsvarer efter de relevante regler og skal behandles som uskyldig, indtil skyld er bevist.

Straffeprocessen skal balancere effektiv opklaring med beskyttelse mod vilkårlige indgreb. Ransagning, varetægtsfængsling, beslaglæggelse og aflytning kræver derfor særlige betingelser, kontrol og proportionalitet.

\section{Påstand, anbringender og beviser}

I en civil sag er \term{påstanden} det resultat, parten beder retten om, eksempelvis betaling af et bestemt beløb. \term{Anbringender} er de juridiske og faktiske grunde, der skal føre til resultatet. \term{Beviser} er materialet, som skal understøtte de faktiske påstande.

En enkel kontrol er:
\begin{itemize}
  \item Hvad skal retten beslutte?
  \item Hvilke betingelser skal være opfyldt?
  \item Hvilket faktum opfylder hver betingelse?
  \item Hvilket bevis viser hvert faktum?
\end{itemize}

Hvis et dokument ikke knytter sig til en betingelse eller et omstridt faktum, er det måske ikke relevant, selv om det føles vigtigt for historien.

\section{Bevisbyrde og bevisvurdering}

Bevisbyrden angiver, hvem der bærer risikoen for, at et faktum ikke bliver lagt til grund. I civile sager ligger bevisbyrden ofte hos den, der påstår et krav eller en undtagelse, men regler og praksis kan flytte eller lette den. I straffesager ligger bevisbyrden for skyld hos anklagemyndigheden.

Bevisvurdering er retten til at vurdere beviserne frit inden for en samlet og begrundet bedømmelse. Fri bevisbedømmelse betyder ikke, at dommeren må vælge efter intuition alene. Begrundelsen skal vise, hvorfor et vidne, et dokument eller en teknisk oplysning er tillagt større vægt end en anden.

\section{Kontradiktion og lighed mellem parter}

Kontradiktion betyder, at en part skal have mulighed for at kende og kommentere det materiale, der kan få betydning for afgørelsen. Det beskytter mod hemmelige argumenter og giver retten et bedre beslutningsgrundlag.

Lighed mellem parterne betyder ikke, at parterne har samme ressourcer eller samme bevisadgang i virkeligheden. Procesregler, vejledning, advokatbeskikkelse og særlige beskyttelsesregler kan derfor være nødvendige for at skabe en fair mulighed for deltagelse.

\section{Dom, appel og fuldbyrdelse}

En dom indeholder resultatet og begrundelsen. Den kan være endelig, eller den kan ankes efter regler om instans, beløbsgrænser og frister. Appel er ikke altid en fuld genbehandling; den konkrete mulighed afhænger af sagstypen og den instans, der har truffet afgørelsen.

Hvis den tabende part ikke frivilligt opfylder dommen, kan vinderen søge fuldbyrdelse. En vundet dom er derfor ikke altid det samme som at have fået pengene eller handlingen i praksis. Solvens, aktiver og fuldbyrdelsesregler kan være afgørende.

\begin{tabularx}{\textwidth}{@{}>{\RaggedRight\arraybackslash}p{2.6cm}X X@{}}
\toprule
\textbf{Spørgsmål} & \textbf{Civil sag} & \textbf{Straffesag} \\
\midrule
Hvem rejser sagen? & Typisk en privat eller offentlig sagsøger. & Staten ved anklagemyndigheden. \\
Formål & Afgøre et krav eller en retlig relation. & Afgøre skyld og eventuel straf. \\
Bevisbyrde & Varierer med krav og regel. & Anklagemyndigheden skal bevise skyld. \\
Resultat & Dom om betaling, pligt, status eller afvisning. & Frifindelse eller dom med straf/sanktion. \\
\bottomrule
\end{tabularx}

\section{Alternativ konfliktløsning}

Forlig, mediation, voldgift og nævnsbehandling kan løse en konflikt uden en fuld retssag. Det kan være hurtigere, billigere og mere fleksibelt, men det kan også begrænse offentlighed, appel eller præcedens. Valget afhænger af kravets art, parternes forhold og den aftale eller lov, der gælder.

\practice{Lav en procesplan med datoer for frister, dokumenter, vidner, klage og appel. Skriv samtidig, hvad der sker, hvis fristen overskrides. Procesretten belønner sjældent en god argumentation, der kommer for sent.}

\question{En person har et veldokumenteret krav, men opdager først efter fristen, at den forkerte virksomhed er sagsøgt. Hvilke procesmæssige problemer kan opstå, og hvorfor er de ikke bare en formalitet?}

\section*{Kapitelopsamling}
\begin{itemize}
  \item Materiel ret og procesret er forskellige, men gensidigt afhængige.
  \item Påstand, anbringender, fakta og beviser skal holdes adskilt.
  \item Bevisbyrden fordeler risikoen ved usikkerhed; fri bevisbedømmelse kræver begrundelse.
  \item Kontradiktion og lighed mellem parter er centrale fair-process-principper.
  \item Dom, appel og fuldbyrdelse er tre forskellige trin i en retssag.
\end{itemize}

